%% ============================================================
%%  Physics-Informed Neural Networks and Quantum Dot Systems
%%  Author: Morten Hjorth-Jensen
%% ============================================================
\documentclass[aspectratio=169,10pt]{beamer}

%% ── Theme & colour ──────────────────────────────────────────────────────────
\usetheme{Madrid}
\usecolortheme{seahorse}
\usefonttheme[onlymath]{serif}

\setbeamercolor{frametitle}{bg=blue!15!white,fg=blue!60!black}
\setbeamercolor{title}{fg=blue!70!black}
\setbeamercolor{block title}{bg=blue!20!white,fg=blue!60!black}
\setbeamercolor{block body}{bg=blue!5!white}
\setbeamercolor{alerted text}{fg=red!60!black}
\setbeamertemplate{navigation symbols}{}

%% ── Packages ────────────────────────────────────────────────────────────────
\usepackage{amsmath,amssymb,bm,physics}
\usepackage{graphicx}
\usepackage{booktabs}
\usepackage{xcolor}
\usepackage{tikz}
\usepackage{makecell}
\usepackage[flushleft]{threeparttable}
\usepackage{url}   % provides \nolinkurl for path display in placeholders

%% ── Figure placeholder macro ─────────────────────────────────────────────────
%%   \figph{path}{caption}  — full-width grey box
%%   \figphc{path}{caption} — compact (side-by-side) box
\newcommand{\figph}[2]{%
  \par\smallskip
  \noindent\colorbox{gray!10}{\parbox{0.88\linewidth}{%
    \centering\smallskip
    \textbf{\color{gray!70}[Figure placeholder]}\\[3pt]%
    {\footnotesize\color{gray!60}\nolinkurl{#1}}\\[3pt]%
    {\small\color{gray!60}#2}\smallskip}}%
  \par\smallskip}
\newcommand{\figphc}[2]{%
  \colorbox{gray!10}{\parbox{0.44\linewidth}{%
    \centering\smallskip
    \textbf{\color{gray!70}[Fig]}\\[2pt]%
    {\tiny\color{gray!60}\nolinkurl{#1}}\\[2pt]%
    {\scriptsize\color{gray!60}#2}\smallskip}}}

%% ── Metadata ────────────────────────────────────────────────────────────────
\title[PINNs and Quantum Dots]{Physics-Informed Neural Networks\\
  and Properties of Quantum Dot Systems}
\author[M.~Hjorth-Jensen]{Morten Hjorth-Jensen}
\institute[UiO]{Department of Physics \& Center for Computing in Science Education\\
  University of Oslo, Norway}
\date{\today}

%% ============================================================
\begin{document}
%% ============================================================

%% ─── Slide 1: Title ─────────────────────────────────────────────────────────
\begin{frame}
  \titlepage
\end{frame}

%% ─── Slide 2: Outline ───────────────────────────────────────────────────────
\begin{frame}{Outline}
  \tableofcontents[hidesubsections]
\end{frame}

%% ============================================================
\section{Motivation and Context}
%% ============================================================

%% ─── Slide 3 ────────────────────────────────────────────────────────────────
\begin{frame}{Quantum Dots as Artificial Atoms}
  \begin{columns}[T]
    \column{0.58\textwidth}
      \begin{itemize}
        \item Semiconductor quantum dots $\equiv$ \alert{artificial atoms}
        \item Precise tunability: particle number $N$, confinement $\omega$,
              interaction strength, external fields
        \item Intersection of condensed matter, quantum information
              and nanotechnology
        \item Rich many-body physics even for few particles:
          \begin{itemize}
            \item Wigner molecule formation
            \item Spin ordering
            \item Non-trivial excitation spectra
          \end{itemize}
        \item Clean realisation of interacting fermions in reduced dimensions:
              Coulomb, exchange, quantum confinement all play essential roles
      \end{itemize}
    \column{0.38\textwidth}
      \begin{block}{Key parameters}
        \begin{tabular}{ll}
          $N$       & particle number \\
          $\omega$  & trap frequency  \\
          $\lambda$ & interaction strength
        \end{tabular}
      \end{block}
      \vspace{0.4cm}
      \begin{alertblock}{Core challenge}
        Hilbert space grows combinatorially with $N$; exact solutions
        rapidly become intractable.
      \end{alertblock}
  \end{columns}
\end{frame}

%% ─── Slide 4 ────────────────────────────────────────────────────────────────
\begin{frame}{Why Existing Methods Struggle}
  \begin{columns}[T]
    \column{0.48\textwidth}
      \textbf{Traditional methods and their limitations}
      \begin{itemize}
        \item \textbf{Exact diagonalisation:} restricted to very small $N$
        \item \textbf{Quantum Monte Carlo:} fermion sign problem;
              stochastic errors
        \item \textbf{DFT:} uncontrolled approximations; limited in
              strongly correlated regimes
        \item \textbf{Mean-field / perturbation:} fails when correlations
              are large
      \end{itemize}
      \vspace{0.3cm}
      All methods struggle when exploring high-dimensional parameter
      spaces or modelling realistic experimental conditions.
    \column{0.48\textwidth}
      \textbf{Machine-learning approaches}
      \begin{enumerate}
        \item Variational Monte Carlo (VMC) with neural ansatz
        \item Autoregressive neural states
        \item \alert{Physics-Informed Neural Networks (PINNs)}
      \end{enumerate}
      \vspace{0.4cm}
      \begin{block}{This talk}
        PINNs within a Slater--Jastrow--neural-correlator ansatz applied
        to quantum dots, $N = 2, 6, 12, 20$, across five orders of
        magnitude in confinement $\omega$.
      \end{block}
  \end{columns}
\end{frame}



%% ─── Slide 5 ────────────────────────────────────────────────────────────────
\begin{frame}{Neural Paradigms for Quantum Many-Body Problems}
  \renewcommand{\arraystretch}{1.35}
  \begin{tabular}{p{2.6cm}p{4.0cm}p{4.6cm}}
    \toprule
    \textbf{Method} & \textbf{Core idea} & \textbf{Natural regime} \\
    \midrule
    Variational MC (VMC) with neural ansatz &
      Minimise $\langle H\rangle$ via MC sampling &
      Large $N$, ground states, lattice systems \\
    Autoregressive neural states &
      Exact unbiased sampling via product of conditionals &
      Discrete degrees of freedom, large scale \\
    \alert{PINNs} &
      Enforce PDE and symmetry residuals directly in loss &
      Continuous space, complex geometries, dynamics \\
    \bottomrule
  \end{tabular}

  \vspace{0.5cm}
  \begin{alertblock}{Key PINN advantage for quantum dots}
    No Monte Carlo sampling; wave function learned as a continuous function
    of particle coordinates; antisymmetry, normalisation, and conservation
    laws imposed as explicit constraints or penalty terms.
  \end{alertblock}
\end{frame}

%% ============================================================
\section{Physics-Informed Neural Networks}
%% ============================================================

%% ─── Slide 6 ────────────────────────────────────────────────────────────────
\begin{frame}{PINN Formulation}
  Let $u:\Omega\subset\mathbb{R}^d\to\mathbb{R}$ satisfy
  $\mathcal{N}[u](x)=0$ on $\Omega$ with $u=g$ on $\partial\Omega$.
  \medskip

  A PINN represents $u$ by a neural network $u_\theta(x)$ and minimises:
  \begin{equation*}
    \mathcal{L}(\theta)
    = \lambda_{\rm data}\!\sum_i|u_\theta(x_i)-y_i|^2
    + \lambda_{\rm PDE}\!\sum_j|\mathcal{N}[u_\theta](\xi_j)|^2
    + \lambda_{\rm BC}\!\sum_k|u_\theta(\zeta_k)-g(\zeta_k)|^2
  \end{equation*}

  \begin{columns}[T]
    \column{0.32\textwidth}
      \begin{block}{Data term}
        Observed data $\{x_i,y_i\}$
      \end{block}
    \column{0.32\textwidth}
      \begin{block}{PDE residual}
        Collocation points $\xi_j$ in $\Omega$
      \end{block}
    \column{0.32\textwidth}
      \begin{block}{BC term}
        Boundary pts $\zeta_k$ on $\partial\Omega$
      \end{block}
  \end{columns}

  \vspace{0.4cm}
  \textbf{Key tool — Automatic differentiation (AD):} evaluates
  $\nabla u_\theta$ and $\Delta u_\theta$ exactly (to floating-point
  precision), enabling end-to-end training constrained by physical laws.
\end{frame}

%% ─── Slide 7 ────────────────────────────────────────────────────────────────
\begin{frame}{Activation Functions and Higher-Order Derivatives}
  \begin{columns}[T]
    \column{0.50\textwidth}
      \begin{itemize}
        \item PINNs require second- (or higher-) order derivatives through AD
        \item Activation function choice is critical for $\ell \ge 2$ PDEs
        \item \textbf{Smooth activations required:} $\tanh$, SiLU, Softplus
        \item \textbf{ReLU is unsuitable} for second-order operators
              (second derivative is zero almost everywhere)
      \end{itemize}
      \vspace{0.3cm}
      \begin{alertblock}{Sobolev approximation}
        Must control error in $W^{\ell,p}$ norms, not just $L^2/L^\infty$.
        Smooth activations ensure well-behaved derivatives.
      \end{alertblock}
    \column{0.46\textwidth}
      \centering
      \figph{figures/activation_derivatives}{Behaviour of activation
        functions and their higher-order derivatives relevant for PINNs.
        (Fig.~in paper.)}
  \end{columns}
\end{frame}

%% ─── Slide 8 ────────────────────────────────────────────────────────────────
\begin{frame}{Five Theoretical Challenges of PINNs}
  \begin{enumerate}\itemsep4pt
    \item \textbf{Sobolev approximation.}
          Universal approximation in $L^2/L^\infty$ does not automatically
          extend to derivatives; need control in $W^{\ell,p}$ norms matching
          the PDE order.

    \item \textbf{Stability (residual-to-error).}
          Small residual $\not\Rightarrow$ small solution error; well-posedness
          is required; inverse problems may need regularisation.

    \item \textbf{Sampling \& generalisation.}
          Collocation minimises a sampled surrogate; bad capacity control or
          imbalanced interior/boundary sampling can leave large population
          residuals.

    \item \textbf{Optimisation and conditioning.}
          Disparate physics terms produce ill-conditioned gradients; NTK
          spectrum governs convergence rate.

    \item \textbf{Spectral bias.}
          Gradient methods prefer low-frequency components; sharp layers or
          oscillations require special treatment (Fourier features, curriculum).
  \end{enumerate}

  \vspace{0.2cm}
  \begin{block}{Reference}
    De Ryck et al.\ (2024) — detailed theoretical treatment.
  \end{block}
\end{frame}

%% ─── Slide 9 ────────────────────────────────────────────────────────────────
\begin{frame}{Mitigation Strategies and Inductive Biases}
  Generalisation in physics ML means \alert{consistency with physical laws},
  not merely accuracy on unseen data. Three interacting mechanisms:
  \begin{enumerate}
    \item Algorithmic stability of gradient-based optimisation
    \item Implicit bias toward smooth, low-complexity solutions
    \item Architectural inductive biases encoding domain knowledge
  \end{enumerate}

  \vspace{0.3cm}
  \begin{block}{Practical engineering checklist}
    \small
    \begin{columns}[T]
      \column{0.48\textwidth}
        \begin{itemize}
          \item Weak/variational or energy formulations when strong form
                is unstable
          \item Hard-enforce essential boundary conditions
          \item Non-dimensionalise and balance loss terms
        \end{itemize}
      \column{0.48\textwidth}
        \begin{itemize}
          \item Smooth activations; architectures controlling derivatives
          \item Exploit symmetry and equivariance
          \item Second-order optimisers (L-BFGS, Stochastic Reconfiguration)
          \item Importance/stratified sampling; space-time curricula
        \end{itemize}
    \end{columns}
  \end{block}
\end{frame}

%% ============================================================
\section{The Ansatz}
%% ============================================================

%% ─── Slide 10 ───────────────────────────────────────────────────────────────
\begin{frame}{Ansatz: Slater--Jastrow--Neural Correlator}
  \begin{equation*}
    \Psi_{\theta,\beta}(R)
    = \underbrace{\mathrm{SD}(R')}_{\text{antisymmetry}}\;
      \exp\!\bigl(
        \underbrace{U(R)}_{\text{analytic pair}}
        +\underbrace{W_\theta(R)}_{\text{neural correlator}}
      \bigr),
    \qquad
    R' = R + \underbrace{\Delta_\beta(R)}_{\text{backflow}}
  \end{equation*}

  \vspace{0.2cm}
  \begin{columns}[T]
    \column{0.52\textwidth}
      \begin{block}{Four components}
        \begin{itemize}
          \item $\mathrm{SD}(R')$ — Slater determinant: exact fermionic
                antisymmetry
          \item $U(R)$ — analytic cusp/pair term (evaluated on \emph{unscaled}
                distances $r_{ij}$)
          \item $W_\theta(R)$ — compact permutation-invariant neural correlator
                (evaluated on \emph{trap-scaled} coords $\tilde{\mathbf x}_i
                = \sqrt\omega\,\mathbf x_i$)
          \item $\Delta_\beta(R)$ — optional backflow, only inside $\mathrm{SD}$
        \end{itemize}
      \end{block}
    \column{0.44\textwidth}
      \begin{alertblock}{Coordinate convention}
        $W_\theta$ uses trap-scaled $\tilde{\mathbf x}_i=\sqrt\omega\,\mathbf x_i$;
        $U$ uses physical $r_{ij}$.\\[5pt]
        This separates the length scales set by the trap from those set by the
        Coulomb interaction.
      \end{alertblock}
      \vspace{0.3cm}
      Two backflow variants:\\
      \textbf{PINN+BF} (\texttt{BackflowNet}) and\\
      \textbf{PINN+CTNN} (\texttt{CTNNBackflowNet})
  \end{columns}
\end{frame}

%% ─── Slide 11 ───────────────────────────────────────────────────────────────
\begin{frame}{Neural Correlator $W_\theta$: Architecture Details}
  \begin{columns}[T]
    \column{0.52\textwidth}
      \textbf{Particle branch (DeepSets)}
      \[
        \mathbf h_i^\phi = \phi(\tilde{\mathbf x}_i),\qquad
        \overline{\mathbf h}^\phi
        = \tfrac{1}{N}\textstyle\sum_i \mathbf h_i^\phi
      \]

      \textbf{Pair branch — safe distance features}\\
      Mollified radius: $\tilde r_{ij}^{\rm soft}=\!\sqrt{\tilde r_{ij}^2
        +\varepsilon^2}$,\; $\varepsilon\propto a_{\rm ho}$

      Six smooth channels per pair:
      \[
        \mathbf u_{ij}=\bigl[s_1,s_2,s_3,{\rm rbf}_{0.25},{\rm rbf}_1,{\rm rbf}_4\bigr]
      \]
      Short-range gate: $\chi(\tilde r)=\tilde r^2/(\tilde r^2+r_g^2)$
      \[
        \overline{\mathbf h}^\psi
        = \tfrac{2}{N(N-1)}\textstyle\sum_{i<j}
          \chi(\tilde r_{ij})\,\mathbf h^\psi_{ij}
      \]
    \column{0.45\textwidth}
      \textbf{Global scalars + readout}
      \[
        g_1 = \tfrac{1}{Nd}\textstyle\sum_i\!\|\tilde{\mathbf x}_i\|^2,\quad
        g_2 = \tfrac{2}{N(N-1)}\textstyle\sum_{i<j} s_1
      \]
      \[
        W_\theta = \rho\!\bigl(\overline{\mathbf h}^\phi\,\|\,
        \overline{\mathbf h}^\psi\,\|\,[g_1,g_2]\bigr)
      \]
      \vspace{0.2cm}
      \begin{alertblock}{Design principles}
        \begin{itemize}
          \item Permutation-invariant by construction
          \item Safe features avoid derivative singularities at $r_{ij}\to0$
          \item Mollification scale $\varepsilon\propto a_{\rm ho}$ tracks the trap
        \end{itemize}
      \end{alertblock}
  \end{columns}
\end{frame}

%% ─── Slide 12 ───────────────────────────────────────────────────────────────
\begin{frame}{Analytic Pair Term and Backflow Modules}
  \begin{columns}[T]
    \column{0.48\textwidth}
      \begin{block}{Analytic cusp pair term $U(R)$}
        \[
          U(R)=\sum_{i<j} u_{\sigma_i\sigma_j}(r_{ij}),\quad
          u_{\sigma\sigma'}(r)=\gamma_{\sigma\sigma'}\,r\,e^{-r}
        \]
        Spin-dependent couplings satisfying cusp conditions:
        \[
          \gamma_{\uparrow\downarrow}=\tfrac{1}{d-1},\qquad
          \gamma_{\uparrow\uparrow}=\tfrac{1}{d+1}
        \]
        Evaluated on \emph{unscaled} distances $r_{ij}=\|\mathbf x_i-\mathbf x_j\|$.
      \end{block}
    \column{0.48\textwidth}
      \begin{block}{Backflow $\Delta_\beta(R)$}
        Shifts only the Slater determinant: $R'=R+\Delta_\beta(R)$.

        \smallskip
        \textbf{\texttt{BackflowNet}}: pairwise message-passing MLP;
        output $s_\beta\tanh(\delta\mathbf x_i)$.

        \smallskip
        \textbf{\texttt{CTNNBackflowNet}}: copresheaf-style module;
        explicit node features $\mathbf h_i\in\mathbb R^{H_v}$ and
        edge features $\mathbf h_{ij}\in\mathbb R^{H_e}$; learned
        linear transports between node and edge spaces.

        \smallskip
        Both: zero-initialised output head; centre-of-mass neutral.
      \end{block}
      \vspace{0.2cm}
      \textbf{Complexity:} dense pairwise $\Rightarrow O(BN^2)$.
  \end{columns}
\end{frame}

%% ============================================================
\section{Results: Ground-State Energies}
%% ============================================================

%% ─── Slide 13 ───────────────────────────────────────────────────────────────
\begin{frame}{Training Pipeline and Benchmark Strategy}
  \begin{columns}[T]
    \column{0.52\textwidth}
      \textbf{Training pipeline}
      \begin{enumerate}
        \item \textbf{Residual-based pretraining:}
              PINN loss enforces the many-body Schr\"odinger equation
              at collocation points
        \item \textbf{Stochastic Reconfiguration (SR) tail:}
              second-order-like variational refinement; improves energy
              and wave-function quality
      \end{enumerate}
      \vspace{0.3cm}
      \textbf{Internal consistency check}\\
      Harmonic-plus-Coulomb virial identity:
      \[
        2\langle T\rangle = 2\langle V_{\rm trap}\rangle
        - \langle V_{\rm int}\rangle
      \]
      Satisfied to sub-percent accuracy (worst case $\lesssim 0.4\%$).
    \column{0.44\textwidth}
      \begin{block}{Benchmark scope}
        \begin{itemize}
          \item $N=2,6,12,20$ electrons
          \item $\omega\in\{10^{-3},10^{-2},0.1,0.5,1.0\}$
          \item Compared to diffusion Monte Carlo (DMC) where available
          \item Two ans\"atze: PINN+BF and PINN+CTNN
        \end{itemize}
      \end{block}
      \vspace{0.3cm}
      \begin{alertblock}{Energy closure}
        $E-(T+V_{\rm int}+V_{\rm trap})\approx 0$\\
        typically $10^{-11}$--$10^{-7}$ (absolute)
      \end{alertblock}
  \end{columns}
\end{frame}

%% ─── Slide 14 ───────────────────────────────────────────────────────────────
\begin{frame}{Ground-State Energies: Comparison with DMC}
  \scriptsize
  \renewcommand{\arraystretch}{1.20}
  \begin{tabular}{c|c|c|cc|cc}
    \toprule
    $N$ & $\omega$ & DMC reference & PINN+BF & \%\,err & PINN+CTNN & \%\,err \\
    \midrule
    2 & 0.01  & $0.073839(2)$  & $0.073838(1)$ & $-0.0014$ & $0.0738382(5)$ & $-0.0014$ \\
      & 1.00  & $3.00000$      & $2.99999(1)$  & $-0.0003$ & $3.00000(2)$   & $-0.0000$ \\
    \midrule
    6 & 0.10  & $3.55385(5)$   & $3.5549(1)$   & $+0.030$  & $\mathbf{3.55388(5)}$ & $+0.0008$ \\
      & 0.50  & $11.78484(6)$  & $11.7895(4)$  & $+0.040$  & $\mathbf{11.7847(2)}$ & $-0.0000$ \\
      & 1.00  & $20.15932(8)$  & $20.1610(6)$  & $+0.008$  & $\mathbf{20.1585(3)}$ & $-0.004$  \\
    \midrule
    12& 0.50  & $39.1596(1)$   & $39.1786(8)$  & $+0.049$  & $\mathbf{39.1604(3)}$ & $+0.002$  \\
      & 1.00  & $65.7001(1)$   & $65.717(1)$   & $+0.026$  & $\mathbf{65.69556(5)}$& $-0.007$  \\
    \midrule
    20& 0.10  & $29.9779(1)$   & ---           & ---       & $29.9888(2)$          & $+0.036$  \\
      & 0.50  & $93.8752(1)$   & ---           & ---       & $\mathbf{93.8789(5)}$ & $+0.004$  \\
      & 1.00  & $155.8822(1)$  & ---           & ---       & $\mathbf{155.8738(7)}$& $-0.002$  \\
    \bottomrule
  \end{tabular}
  \normalsize

  \vspace{0.25cm}
  \begin{itemize}
    \item $N=2$: both ans\"atze reproduce DMC within statistical error for
          all $\omega$
    \item $N\ge6$: PINN+CTNN systematically more accurate than PINN+BF
    \item $N=20$: CTNN remains accurate at scale; errors $\lesssim 0.04\%$
  \end{itemize}
\end{frame}


\begin{frame}[allowframebreaks]{Full Ground-State Energy Table}
  \scriptsize
  \renewcommand{\arraystretch}{1.15}
  \begin{tabular}{c|c|c|cc|cc}
    \toprule
    $N$ & $\omega$ & DMC ref. & PINN+BF & \%\,err & PINN+CTNN & \%\,err \\
    \midrule
    2 & $10^{-3}$ & ---            & $0.0137948(8)$ & ---        & $0.013778(1)$    & ---      \\
      & $10^{-2}$ & $0.073839(2)$  & $0.073838(1)$  & $-0.0014$ & $0.0738382(5)$   & $-0.0014$\\
      & $0.10$    & $0.44079(1)$   & $0.44079(1)$   & $+0.0000$ & $0.440796(4)$    & $-0.0000$\\
      & $0.50$    & $1.65977(1)$   & $1.65976(2)$   & $-0.0006$ & $1.65976(1)$     & $-0.0006$\\
      & $1.00$    & $3.00000$      & $2.99999(1)$   & $-0.0003$ & $3.00000(2)$     & $-0.0000$\\
    \midrule
    6 & $10^{-3}$ & ---            & $0.140913(1)$  & ---        & $0.1408172(7)$   & ---      \\
      & $10^{-2}$ & ---            & $0.69125(2)$   & ---        & $0.69036(1)$     & ---      \\
      & $0.10$    & $3.55385(5)$   & $3.5549(1)$    & $+0.0295$ & $3.55388(5)$     & $+0.0008$\\
      & $0.50$    & $11.78484(6)$  & $11.7895(4)$   & $+0.0395$ & $11.7847(2)$     & $-0.0000$\\
      & $1.00$    & $20.15932(8)$  & $20.1610(6)$   & $+0.0083$ & $20.1585(3)$     & $-0.0041$\\
    \midrule
    12& $10^{-3}$ & ---            & $0.515823(3)$  & ---        & $0.515365(4)$    & ---      \\
      & $10^{-2}$ & ---            & $2.48620(5)$   & ---        & $2.47363(4)$     & ---      \\
      & $0.10$    & $12.26984(8)$  & $12.2731(2)$   & $+0.0160$ & $12.2718(1)$     & ---      \\
      & $0.50$    & $39.1596(1)$   & $39.1786(8)$   & $+0.0485$ & $39.1604(3)$     & $+0.0018$\\
      & $1.00$    & $65.7001(1)$   & $65.717(1)$    & $+0.0257$ & $65.69556(5)$    & $-0.0069$\\
    \midrule
    20& $10^{-3}$ & ---            & ---            & ---        & $1.293033(6)$    & ---      \\
      & $10^{-2}$ & ---            & ---            & ---        & $6.14645(5)$     & ---      \\
      & $0.10$    & $29.9779(1)$   & ---            & ---        & $29.9888(2)$     & $+0.0364$\\
      & $0.50$    & $93.8752(1)$   & ---            & ---        & $93.8789(5)$     & $+0.0040$\\
      & $1.00$    & $155.8822(1)$  & ---            & ---        & $155.8738(7)$    & $-0.0024$\\
    \bottomrule
  \end{tabular}
\end{frame}



%% ─── Slide 15 ───────────────────────────────────────────────────────────────
\begin{frame}{Energy Partitioning: Approaching the Wigner Regime}
  \begin{columns}[T]
    \column{0.52\textwidth}
      Total energy decomposition:
      \[
        E = \langle T\rangle
            + \langle V_{\rm int}\rangle
            + \langle V_{\rm trap}\rangle
      \]
      Two diagnostics tracking interaction dominance:
      \[
        \Gamma(\omega)
        = \frac{\langle V_{\rm int}\rangle}{\langle T\rangle},\qquad
        \mathcal{C}(\omega)
        = \frac{2\langle V_{\rm trap}\rangle}{\langle V_{\rm int}\rangle}
      \]
      \begin{itemize}
        \item $\Gamma\!\uparrow$ as $\omega\!\downarrow$: interactions dominate
        \item $\mathcal{C}\to 1$: quasi-classical trap--Coulomb balance
      \end{itemize}
    \column{0.44\textwidth}
      \scriptsize
      \renewcommand{\arraystretch}{1.2}
      \begin{tabular}{c|cc|cc}
        \toprule
        $N$ & $\Gamma_{\omega=1}$ & $\mathcal{C}_{\omega=1}$
            & $\Gamma_{\omega=10^{-3}}$ & $\mathcal{C}_{\omega=10^{-3}}$ \\
        \midrule
        2  & 0.93 & 3.14 & 8.64  & 1.226 \\
        6  & 2.41 & 1.83 & 24.55 & 1.074 \\
        12 & 3.61 & 1.56 & 41.11 & 1.049 \\
        20 & 4.66 & 1.43 & 57.26 & 1.038 \\
        \bottomrule
      \end{tabular}
      \normalsize
      \vspace{0.3cm}
      \begin{alertblock}{Key trend}
        At $\omega=10^{-3}$, $N\ge 12$: $\mathcal{C}< 1.05$, i.e.\
        $V_{\rm int}\approx 2V_{\rm trap}$ to $\lesssim 5\%$.
        System is well into interaction-dominated regime.
      \end{alertblock}
  \end{columns}
\end{frame}

%% ─── Slide 16 ───────────────────────────────────────────────────────────────
\begin{frame}{Energy-Partition Diagnostics (Figures)}
  \begin{columns}[T]
    \column{0.49\textwidth}
      \centering
      \figphc{figures/energy/ratios_allN_Gamma.png}{%
        (a) $\Gamma(\omega)=\langle V_{\rm int}\rangle/\langle T\rangle$
        (log--log). All $N$. Interaction dominance rises sharply as
        $\omega\to 0$. (Fig.~\ref{fig:energy_ratios}a)}
    \column{0.49\textwidth}
      \centering
      \figphc{figures/energy/ratios_allN_classical.png}{%
        (b) $\mathcal{C}(\omega)=2\langle V_{\rm trap}\rangle/
        \langle V_{\rm int}\rangle$. Horizontal line: $\mathcal{C}=1$.
        Larger $N$ approaches the quasi-classical limit more rapidly.
        (Fig.~\ref{fig:energy_ratios}b)}
  \end{columns}

  \vspace{0.3cm}
  RMS cloud radius per particle:
  $r_{\rm rms}=\sqrt{2\langle V_{\rm trap}\rangle/(N\omega^2)}$

  \begin{center}
    \begin{tabular}{c|cccc}
      \toprule
      $\omega$ & $N=2$ & $N=6$ & $N=12$ & $N=20$ \\
      \midrule
      $1$         &  1.14 &  1.62 &  2.04 &  2.40 \\
      $10^{-3}$   & 69.9  & 126   & 171   & 209   \\
      \bottomrule
    \end{tabular}
  \end{center}
  \small Dramatic swelling confirms interaction-driven expansion.
\end{frame}


%% ============================================================
\section{Summary and Outlook}
%% ============================================================

%% ─── Slide 27 ───────────────────────────────────────────────────────────────
\begin{frame}{Summary: Energetics}
  \begin{columns}[T]
    \column{0.50\textwidth}
      \begin{block}{Energetic accuracy}
        \begin{itemize}
          \item Residual pretraining + SR tail pipeline
          \item Virial identity: internal validation to $<0.4\%$
          \item $N=2$: both PINN+BF and PINN+CTNN agree with DMC
                to within statistical uncertainty
          \item $N=6,12,20$: PINN+CTNN systematically outperforms
                PINN+BF at moderate/strong confinement
          \item Deviations from DMC: $\lesssim 0.05\%$ in benchmarked
                regimes
        \end{itemize}
      \end{block}
    \column{0.50\textwidth}
      \begin{block}{Energy partitioning}
        \begin{itemize}
          \item $\Gamma(\omega)$: interaction-to-kinetic ratio grows
                from $<1$ to $\sim 57$ across the $\omega$ range
          \item $\mathcal{C}(\omega)\to 1$: quasi-classical balance
                at $\omega=10^{-3}$; better approached for larger $N$
          \item $r_{\rm rms}$: $\sim\!70\times$ expansion from
                $\omega=1$ to $\omega=10^{-3}$
        \end{itemize}
      \end{block}
  \end{columns}
\end{frame}


%% ─── Slide 29 ───────────────────────────────────────────────────────────────
\begin{frame}{Methodological Contributions}
  \begin{enumerate}\itemsep5pt
    \item \textbf{Architecture.}
          Permutation-invariant DeepSets + safe pair features; mollified
          distances avoid coalescence singularities; PINN loss with explicit
          physical constraints.

    \item \textbf{Training pipeline.}
          Residual-based pretraining $\to$ SR tail; virial identity used as
          ongoing physics-based validation.

  \end{enumerate}

  \vspace{0.3cm}
  \begin{alertblock}{Broader applicability}
    Same PINN framework is directly applicable to:
    other 2D/3D trapped systems $\cdot$ real-time dynamics $\cdot$
    excited states $\cdot$ open quantum systems $\cdot$
    device-level quantum dot modelling.
  \end{alertblock}
\end{frame}

%% ─── Slide 30 ───────────────────────────────────────────────────────────────
\begin{frame}{PINNs vs.\ VMC and Autoregressive States}
  \begin{block}{Why PINNs generalise for physics}
    Inductive biases are not optional — they define what ``correct'' means:
    symmetry, conservation laws, and PDE structure are built into the loss.
  \end{block}

  \vspace{0.3cm}
  \begin{columns}[T]
    \column{0.48\textwidth}
      \textbf{PINNs vs.\ VMC}
      \begin{itemize}
        \item No MC sampling in configuration space
        \item Wave function is a continuous function of coordinates
        \item Natural access to excited states, complex geometries,
              time-dependent problems
        \item Antisymmetry and normalisation as modular constraints
      \end{itemize}
    \column{0.48\textwidth}
      \textbf{PINNs vs.\ autoregressive states}
      \begin{itemize}
        \item No discretisation of continuous space required
        \item Constraints explicit in loss; physics transparent
        \item Better interpretability for analysis of structure
        \item Weaker for very large $N$; complementary, not competing
      \end{itemize}
  \end{columns}

  \vspace{0.3cm}
  \textbf{Conclusion:} VMC excels at large $N$ ground states;
  PINNs may be good candidates for continuous-space, dynamics, constrained, and
  few-to-intermediate-body problems.  Cross-fertilisation is the key
  direction.
\end{frame}

%% ─── Slide 31 ───────────────────────────────────────────────────────────────
\begin{frame}{Open Questions and Future Directions}
  \begin{columns}[T]
    \column{0.50\textwidth}
      \textbf{Computational / methodological}
      \begin{itemize}
        \item Scaling to $N>20$: sparse/message-passing architectures
              to reduce $O(N^2)$ cost
        \item Transfer learning: pre-train on one $\omega$, fine-tune
              across the full $\omega$ grid
        \item Excited states: add orthogonality constraints to loss
        \item Real-time dynamics: extend PINN loss with time derivative
        \item Improved optimisers: natural gradient, KFAC for SR tail
      \end{itemize}
    \column{0.50\textwidth}
      \textbf{Physics questions}
      \begin{itemize}
        \item Detailed melting scenario at intermediate $\omega$
        \item Spin degree of freedom and magnetic field effects
        \item Coupled dots and transport / conductance calculations
        \item Entanglement structure of the many-body wave function
        \item Connection to 2D Wigner crystal in the thermodynamic limit
      \end{itemize}
  \end{columns}
\end{frame}

%% ─── Slide 32 ───────────────────────────────────────────────────────────────
\begin{frame}{Conclusions}
  \begin{enumerate}\itemsep6pt
    \item \textbf{PINNs are accurate} for quantum dots at $N=2$--$20$
          across five orders of magnitude in $\omega$, matching DMC
          benchmarks to $\lesssim 0.05\%$.


    \item \textbf{Energy partitioning and virial validation} confirm
          physical correctness even without DMC references.


    \item \textbf{PINNs provide a unified, physics-grounded framework}
          for few-to-intermediate quantum many-body problems in
          continuous space — with natural extensions to dynamics and
          excited states.
  \end{enumerate}
\end{frame}





\end{document}
